वाक्यसंच एका अर्थाने त्या रचनांचे वर्णन करतात ज्यांमध्ये ती वाक्ये एकत्रितपणे सत्य असतात. त्या रचना म्हणजे त्या वाक्यसंचांची \emph{प्रतिमाने}. उलट, एखाद्या रचनेपासून किंवा रचनांच्या संचापासून सुरुवात केली, तर ती रचना किंवा त्या रचना ज्या वाक्यांची प्रतिमाने आहेत त्या वाक्यांचा संच विचारात घेता येतो. ती त्या रचनेची किंवा रचनांच्या संचाची \emph{उपपत्ती} होय. अशा प्रत्येक वाक्यसंचातून तार्किकरीत्या निष्पन्न होणारे प्रत्येक वाक्य त्या संचात आधीच असते. म्हणजेच तो संच \emph{बंद} असतो. अधिक सर्वसाधारणपणे, संच $\Gamma$ तार्किक निष्पन्नतेखाली बंद असेल, तर त्याला उपपत्ती म्हणतो. $\Delta$ पासून तार्किकरीत्या निष्पन्न होणाऱ्या सर्व वाक्यांचा संच म्हणजे $\Gamma$ असेल, तर $\Gamma$ ही उपपत्ती~$\Delta$ \emph{द्वारे स्वयंसिद्धकीय रीतीने निरूपित} होते असे म्हणतो.

गणितात उपपत्तींची अनेक उदाहरणे मिळतात. उदाहरणार्थ, रेषीय क्रमांची उपपत्ती, गटांची उपपत्ती आणि अंकगणिताच्या उपपत्ती, जसे पेआनो यांच्या स्वयंसिद्धकांद्वारे निरूपित उपपत्ती. इतर विद्याशाखांमध्येही महत्त्वाच्या उपपत्तींची पुष्कळ उदाहरणे आहेत. उदाहरणार्थ, संबंधाधारित डेटाबेस उपपत्ती म्हणून पाहता येतात, आणि सत्तामीमांसेत स्वयंसिद्धकांद्वारे निरूपित करता येतील अशा अवयवत्वाच्या उपपत्तींचा विचार होतो.

अभ्यासासाठी उपपत्ती मांडताना तिच्या साहाय्याने ज्या सर्व गोष्टींविषयी बोलायचे आहे त्या मांडण्याइतकी तिची भाषा अभिव्यक्तिक्षम आहे का, हा एक महत्त्वाचा प्रश्न असतो. ज्या गोष्टी सिद्ध करायच्या आहेत त्या सिद्ध करण्याइतकी ती उपपत्ती समर्थ आहे का, हा दुसरा प्रश्न. संबंध \emph{व्यक्त} करण्यासाठी आवश्यक तितके मुक्त चर असलेले सूत्र लागते. \emph{संचसिद्धान्तात} $\in$ हे एकमेव संबंधचिन्ह आहे, तरी त्याच्या साहाय्याने $\lforall[u][(u \in x \lif u \in y)]$ वापरून $x \subseteq y$ व्यक्त करता येते. मोजकी स्वयंसिद्धके आणि $\subseteq$ च्या व्याख्येसारख्या, क्रमाक्रमाने अधिक गुंतागुंतीच्या होत जाणाऱ्या व्याख्यांची साखळी यांच्या साहाय्याने \emph{निवडीच्या स्वयंसिद्धकासह झर्मेलो--फ्रेंकेल संच उपपत्ती}~$\Log{ZFC}$ (जवळजवळ) प्रत्येक गणितीय दावा व्यक्त करण्याइतकी आणि (जवळजवळ) प्रत्येक गणितीय प्रमेय सिद्ध करण्याइतकी समर्थ आहे.
